\section{The original logarithmic connection and its synchronized support decomposition}
\label{sec:heat-lattice-receiver}

Use the complete convergent substitution proved in the preceding
integral-lattice derivation: $\xi=F(t,y)=y+tB(t,y)$, from
\[
 A=\mathcal O[\xi]/W(t,\xi)
       \xrightarrow[\xi\mapsto F(t,y)]{\ \sim\ }
 A_H=\mathcal O[y]/P_m(t,y),\qquad
 P_m(t,y)=\prod_{a=1}^r(y^2-b_a^2t)\,y^\varepsilon,
 \tag{HL1}
\]
where $\mathcal O=\mathbb C\{t\}$, $r=\lfloor m/2\rfloor$,
$\varepsilon=m-2r$, and $b_a$ are the positive roots of HF2.
The full actual coefficients remain in $F$ and its inverse.
Let $T$ take a coefficient column in the original $\xi$ basis to
its coefficient column in the $y$ basis. The proved substitution
has $T\in\mathrm{GL}_m(\mathcal O)$ and $T(0)=I$.
The calculations below give the exact receivers of this isomorphism;
they do not substitute $P_m$ for the original arithmetic function.

\paragraph{HL1. The connection on the original coefficient lattice.}
Weighted homogeneity gives
$2t\partial_tP_m+y\partial_yP_m=mP_m$.
Consequently on the generic fibre the original root connection,
whose horizontal coordinates are branch evaluations, is
$\partial_t+(y/(2t))\partial_y$ in $A_H$.
On its ordered power basis it has matrix
$D/(2t)$, where $D=\operatorname{diag}(0,1,\ldots,m-1)$.
The chain rule for the exact coefficient map $T$ therefore proves
\[
 \Gamma_A=T^{-1}\left(\frac{D}{2t}T+T'\right),
 \qquad t\nabla A\subseteq A,
 \qquad \operatorname{Res}_{t=0}\nabla\big|_A=\frac D2.
 \tag{HL2}
\]
Indeed for an original column $f$, the transformed covariant
derivative is $(Tf)'+(D/(2t))Tf$; multiplying by $T^{-1}$ gives
HL2. Since $T$ and its inverse are analytic and $T(0)=I$, the
matrix has only a simple pole and the stated residue.
This is the actual original connection, including the unit term
$2U_\xi/U$ proved in HH8; the equality follows from their identical
branch-evaluation derivatives under HL1.

The original monodromy is thus the regular algebra involution
\[
 M_A=T^{-1}\operatorname{diag}((-1)^k)_{k=0}^{m-1}T,
       \qquad M_A(A)=A.
 \tag{HL3}
\]
It is the conjugate of $y\mapsto-y$, so it squares to the identity
and preserves the unit and product. At the endpoint it sends
$\xi^k\mapsto(-1)^k\xi^k$ with every nilpotent retained.

\paragraph{HL2. The conductor index is the exact residue shift.}
Let $B$ be the integral closure of HC5, and use its ordered
component basis $1,u_a$ with $u_a^2=t$, and the central $1$ when
$\varepsilon=1$. Its connection residue has $r+\varepsilon$
zeros and $r$ copies of $1/2$. The original inclusion is the full
matrix $J=C(t)\operatorname{diag}(t^{\lfloor k/2\rfloor})$ of HC7,
with $C(t)$ analytic and invertible. The same chain rule gives
\[
 \Gamma_A=J^{-1}(\Gamma_BJ+J'),\qquad
 \operatorname{Tr}\Gamma_A-\operatorname{Tr}\Gamma_B
    =\frac d{dt}\log\det J
    =\frac{\delta}{t}+\frac d{dt}\log\det C(t),
 \tag{HL4}
\]
where
\[
 \delta=\sum_{k=0}^{m-1}\lfloor k/2\rfloor
   =\left\lfloor\frac{(m-1)^2}{4}\right\rfloor
   =\operatorname{length}_{\mathcal O}(B/A).
 \tag{HL5}
\]
The logarithmic derivatives in HL4 denote $f'/f$ for the displayed
nonzero determinants, so they require no choice of logarithm.
The trace residue difference also equals
$\sum_{k=0}^{m-1}k/2-r/2=m(m-1)/4-r/2=\delta$.
Thus the full coefficient lattice and the integral closure differ
by the precisely evaluated integer parts of their connection
exponents. The analytic determinant unit is retained in HL4.

\paragraph{HL3. Fixed and anti-fixed amplitudes.}
Define $Q(t,z)=\prod_{a=1}^r(z-b_a^2t)$, with empty product $1$.
The coefficient parity of $P_m(t,y)$ proves the following exact
module decomposition under HL1:
\[
 \begin{array}{c|c|c}
 & A^+=\ker(M_A-I)&A^-=\ker(M_A+I)\\ \hline
 m=2r&\mathcal O[z]/Q(t,z)&y\,\mathcal O[z]/Q(t,z)\\
 m=2r+1&\mathcal O[z]/(zQ(t,z))&y\,\mathcal O[z]/Q(t,z).
 \end{array}
 \tag{HL6}
\]
Here $z=y^2$ and the notation in the two columns is carried back
to the original coordinate by the inverse of HL1. The left column
is a subalgebra. The right column is a module, with its products
landing in the left column; it is not asserted to be a subalgebra.
For even $m$, the relation is exactly $Q(t,y^2)=0$; separating
even and odd coefficients gives both columns. For odd $m$, the
relation $yQ(t,y^2)=0$ gives $zQ(t,z)=0$ in even degree and
$Q(t,z)=0$ in odd degree. The bases $1,y,\ldots,y^{m-1}$ show
that there are no further relations. This proves HL6 including
the empty odd column at $m=1$.

For real $t$, the substitution coefficients are real, by uniqueness
of interpolation from the real Taylor coefficients of $H_t$.
Thus $M_A$ commutes with reflection and preserves $T_t$.
If $a\in A^+$, $b\in A^-$, then
$T_t(a,b)=T_t(M_Aa,M_Ab)=-T_t(a,b)$, so
\[
 T_t(A^+,A^-)=0.
 \tag{HL7}
\]
This holds as an identity of the full analytic Gram matrices,
including the endpoint. The endpoint socle $\mathbb C\xi^{m-1}$
lies in the plus column for odd $m$ and in the minus column for
even $m$. Its deformed trace sign from HF10 at negative $t$ is
therefore exactly its monodromy sign. Its actual endpoint trace
is still zero when $m\ge2$.

\paragraph{HL4. The supported projections have an exact common residue.}
The amplitude projections $E_\pm=(I\pm M_A)/2$ obey
$E_+E_-=0$ and $E_++E_-=I$. Lift them to the full supported
$G_L(\mathcal O)$-semimodule $G_L(A)$ by
\[
 \widehat E_\pm(a,\lambda)=(E_\pm a,\lambda),
       \qquad \kappa(a,\lambda)=(0,\lambda).
 \tag{HL8}
\]
Only zero amplitude is allowed below the top support, as in the
original definition. These formulas are consequently well-defined
and additive, and they commute with every scalar in
$G_L(\mathcal O)$: amplitude linearity and the meet of supports
give the assertion directly. Exact composition now gives
\[
 \widehat E_\pm^2=\widehat E_\pm,\qquad
 \widehat E_++\widehat E_-=I,\qquad
 \boxed{\widehat E_+\widehat E_-
            =\widehat E_-\widehat E_+=\kappa.}
 \tag{HL9}
\]
The last map preserves the original support while its amplitude
vanishes. For top-supported input it returns $e$, and for a zero
label $z_\lambda$ it returns that same label. It returns $\tau$
only at $\tau$. Thus no disappearance of a supported component
is inferred from the ordinary amplitude product of projections.

The exact inverse construction is the synchronized module
\[
 \mathcal S=\{((a_+,\lambda),(a_-,\lambda)):
       a_\pm\in A^\pm,\text{ with the allowed original supports}\}.
 \tag{HL10}
\]
The map $G_L(A)\to\mathcal S$ is
$(a,\lambda)\mapsto((E_+a,\lambda),(E_-a,\lambda))$;
its inverse is $((a_+,\lambda),(a_-,\lambda))\mapsto
(a_++a_-,\lambda)$. Addition uses joins of the common supports.
Equip $\mathcal S$ with multiplication
\[
 ((a_+,a_-),\lambda)((b_+,b_-),\mu)
   =((a_+b_++a_-b_-,\ a_+b_-+a_-b_+),\lambda\wedge\mu).
 \tag{HL11}
\]
The parity product rules from HL6 show that both components land
in the stated modules. Expansion of $(a_++a_-)(b_++b_-)$ proves
that the displayed inverse preserves multiplication and the unit.
It is therefore an exact semiring isomorphism for HL11, in
addition to its semimodule isomorphism. The whole lattice $L$
is shared, never replaced by two independent copies or by a
Boolean scalar. Both the amplitude map and the support map commute
with this isomorphism. This is the supported receiver of the
regular original collision monodromy.

\paragraph{HL5. The positive eight-state collision algebra and its full support.}
The same proved $C_2$-projection construction has a concrete receiver
for the new eight-state intake. Use its original critical fibre,
with all coordinates retained:
\[
 \mathcal A_c=\mathcal L\times K_-\times K_\infty,\qquad
 \mathcal L=\mathbb C[T]/T^4,\quad
 K_-=\mathbb C[S]/(S^2-9),\quad
 K_\infty=\mathbb C[\Theta]/(\Theta^2+1).
 \tag{HL12}
\]
The original collision coordinate is $\epsilon=T^2/6$.
These are the original chart and collision formulas of
\href{https://github.com/KokunoYumeto/zeta-function-research-reader/blob/64a77de144720ea9e07943d6ed6f63a178d462cf/workbenches/splitzero-tandem/branches/tau-zero-prime-positivity/EIGHT_STATE_HEAT_COMPARISON.md}{\emph{The signed eight-state completion and its exact heat receivers}}.
The computations required here are given in full below.
The fibre involution used here conjugates coefficients and sends
$T,S,\Theta$ to their negatives. It is the conjugate-deck involution.
The native target-reflection involution instead maps this target
to its reflected target, as proved by that source's equations
ESH17--18; no self-involution of a single moving fibre is inferred
from that separate map.

Writing $x=(\ell,c+dS,\eta_0+\eta_1\Theta)$, multiplication
in the three unchanged power bases gives
\[
 Q_c(x)=4|\ell(0)|^2+2|c|^2-18|d|^2
                         +2|\eta_0|^2+2|\eta_1|^2.
 \tag{HL13}
\]
Indeed the trace on $\mathcal L$ is four times the constant
coefficient; every higher multiplication power of $T$ has zero
diagonal. In $K_-$, $S^\#S=-9$ and the trace of $1$ is two.
In $K_\infty$, $\Theta^\#\Theta=1$ and the same trace is two.
All cross terms have zero scalar coefficient. This proves the
whole sesquilinear form by polarization, with inertia $(4,1,3)$.

Define the actual algebra involution
$\sigma_c(\ell,c+dS,z)=(\ell,c-dS,z)$.
Its fixed subalgebra and its anti-fixed module are
\[
 \mathcal A_c^+=\mathcal L\times\mathbb C\times K_\infty,
 \qquad \mathcal A_c^-=\{(0,dS,0):d\in\mathbb C\}.
 \tag{HL14}
\]
It commutes with the displayed conjugate-deck involution.
The group-algebra homomorphism
$\mathbb C[C_2]\to\operatorname{End}_{\mathbb C}(\mathcal A_c)$,
whose generator maps to $\sigma_c$, sends $(1\pm g)/2$ to
the two projections in HL14. The analogous homomorphism for
the original heat algebra sends $g$ to $M_A$. These give the
exact common domain and the two concrete receivers of HL8--11;
they do not identify their distinct coefficient algebras.

The fixed projection $E_c=(I+\sigma_c)/2$ satisfies the exact
identities
\[
 \begin{gathered}
 Q_c(x)=Q_c(E_cx)-18|d|^2,\qquad
 \operatorname{inertia}(Q_c|_{\mathcal A_c^+})=(4,0,3),\\
 E_c(xy)-E_c(x)E_c(y)
      =(0,9d_xd_y,0).
 \end{gathered}
 \tag{HL15}
\]
The product formula follows by expanding
$(c_x+d_xS)(c_y+d_yS)$ and retaining $S^2=9$.
It evaluates the complete multiplicative defect of this particular
positive projection. The length-four local algebra is unchanged.
Its relative trace to $\mathbb C[\epsilon]/\epsilon^2$ still sends
$T^\#T=-6\epsilon$ to $-12\epsilon$, whereas the scalar trace
sends it to zero. Both maps and the nonzero relative value follow
from the free relative basis $(1,T)$.

The fixed algebra has dimension seven. A nonnegative subspace
for an inertia-$(4,1,3)$ form has dimension at most seven:
modulo its three-dimensional radical the projection of a
nonnegative subspace to the positive four-space is injective,
since its kernel would be a nonzero negative vector.
Thus this fixed subalgebra reaches that bound. It is the unique
largest nonnegative subalgebra containing the complete local and
infinity factors and their central units: in the remaining
$K_-\simeq\mathbb C^2$, a unital subalgebra containing unequal
coordinates contains both point idempotents by linear interpolation,
and then contains the whole indefinite factor. Its only
nonnegative unital alternative is the diagonal copy of $\mathbb C$.

Applying HL8--11 to $\sigma_c$ proves the full supported
projection formula, with the same $\kappa$ and every original
label. The inclusion $G_L(\mathcal A_c^+)\hookrightarrow
G_L(\mathcal A_c)$ is a unital semiring map. On arithmetic
prime spectra its contravariant map identifies the two reduced
$S=\pm3$ points and preserves the local collision point and
both infinity points. This follows from the factorwise
inclusions $\mathbb C\hookrightarrow\mathbb C^2$ and the two
identity maps in HL14. On support primes it is the identity:
an ideal determined by a lattice prime pulls back to the same
lattice prime, since every label is fixed. Hence the exact
full support spectrum map is also evaluated. This positive
subalgebra changes the two reduced $S$ directions and retains
the complete infinitesimal and its relative trace.
